\LessonHeader{9}{Three Doors, One Car and a Puzzle That Fools Everyone}{三扇门、一辆车和一道人人上当的难题}{专题报道 · Feature}{话题\ 概率谜题·思维误区\enspace|\enspace 新词\ \textbf{53}\enspace|\enspace 新义\ \textbf{37}\enspace|\enspace 短语\ \textbf{21}\enspace|\enspace 全文\ \textbf{836}\ 词}{Lesson 9\enspace·\enspace Three Doors, One Car and a Puzzle That…}
\begin{lessonbody}
\PARA{1}I first \NS{}{met}{}{新义}{meet}{/miːt/}{vt.}{遇到；碰到} this \GL{}{puzzle}{}{/ˈpʌzl/}{n.[C]}{难题；谜题} in a \RV{}{high}{} school \NW{}{classroom}{.}{/ˈklɑːsruːm/}{n.[C]}{教室} \SMK{y}The \NW{}{subject}{}{/ˈsʌbdʒɪkt/}{n.[C]}{学科；科目} was \GL{}{statistics}{,}{/stəˈtɪstɪks/}{n.[U]}{统计学} and my \NW{}{teacher}{}{/ˈtiːtʃə(r)/}{n.[C]}{老师} \NWB{}{presented}{}{present}{/prɪˈzent/}{vt.}{提出；展示} a \NS{}{problem}{}{新义}{}{/ˈprɒbləm/}{n.[C]}{（数学）题；难题} that has \PH{\RVB{}{stayed}{}{stay}\ with}{stay with sb}{萦绕心头} me \RV{}{ever}{} \RV{}{since}{.}\EMK{y}{1} 
\YI{1}{那门课是统计学，老师提出了一道题，从那以后它一直萦绕在我心头。}
It is named after \PN{}{Monty Hall}{,}{/ˌmɒnti ˈhɔːl/}{n.}{蒙蒂·霍尔（主持人）} \NS{}{host}{}{新义}{}{/həʊst/}{n.[C]}{主持人} of the \RV{}{TV}{} \NS{}{show}{}{新词性}{}{/ʃəʊ/}{n.[C]}{节目} \PN{}{Let’s Make a Deal}{.}{/ˌlets ˌmeɪk ə ˈdiːl/}{n.}{《我们做笔交易吧》（节目名）}

\PARA{2}There are three \NWB{}{doors}{}{door}{/dɔː(r)/}{n.[C]}{门} in \RV{}{front}{} of you. A new \RV{}{car}{} is behind one of them, and \GL{}{goats}{}{/ɡəʊts/}{n.[C]}{山羊} are behind the other two. You \NS{}{pick}{}{新义}{}{/pɪk/}{vt.}{挑选} \AG{}{door}{} 1. The \AG{}{host}{,} who knows where the \AG{}{car}{} is, opens \RV{}{another}{} \AG{}{door}{} and \NS{}{shows}{}{新义}{show}{/ʃəʊ/}{vt.}{给……看} you a goat. Then he asks if you want to \RV{}{keep}{} \AG{}{door}{} 1 or \GL{}{switch}{}{/swɪtʃ/}{vi.}{换；调换} to the other \NW{}{closed}{}{/kləʊzd/}{adj.}{关着的} \AG{}{door}{.}

\PARA{3}\RV{}{Most}{} people think it is \PH{50-50.}{50-50}{五五开；各占一半} But you \NW{}{should}{}{/ʃʊd/}{modal v.}{应该} \RV{}{always}{} switch: you \RV{}{win}{} \PH{two\ times\ out\ of\ three}{… out of …}{……中有……} if you switch and one time \RV{}{in}{} three if you \RV{}{stay}{.} Switching \NS{}{doubles}{}{新词性}{double}{/ˈdʌbl/}{vt.}{使加倍} your \NS{}{chance}{}{新词性}{}{/tʃɑːns/}{n.[C]}{可能性} of \AGB{}{winning}{.}{win}

\PARA{4}\RVB{}{Students}{}{student} all over the \RV{}{world}{} \RV{}{learn}{} about it \PH{in\ \NS{}{class}{}{新义}{}{/klɑːs/}{n.[U]}{上课；课}}{in class}{在课堂上} \RV{}{during}{} \RV{}{maths}{} and statistics \NWB{}{lessons}{.}{lesson}{/ˈlesn/}{n.[C]}{课} It \RVB{}{became}{}{become} \RV{}{famous}{} in 1990, \RV{}{when}{} \RV{}{writer}{} \PN{}{Marilyn vos Savant}{}{/ˌmærəlɪn vɒs səˈvɑːnt/}{n.}{玛丽莲·沃斯·莎凡特} \NWB{}{explained}{}{explain}{/ɪkˈspleɪn/}{vt.}{解释} it in a \NW{}{magazine}{.}{/ˌmæɡəˈziːn/}{n.[C]}{杂志} About 10,000 \RVB{}{readers}{}{reader} who \NWB{}{disagreed}{}{disagree}{/ˌdɪsəˈɡriː/}{vi.}{不同意} with her wrote \NW{}{angry}{}{/ˈæŋɡri/}{adj.}{生气的} \NWB{}{letters}{,}{letter}{/ˈletə(r)/}{n.[C]}{信} and \RV{}{around}{} 1,000 of them had a \NS{}{doctor}{’s}{新义}{}{/ˈdɒktə(r)/}{n.[C]}{博士} \NW{}{degree}{.}{/dɪˈɡriː/}{n.[C]}{学位}

\PARA{5}\PH{\RV{}{Pretty}{}\ much}{pretty much}{几乎} \NW{}{everybody}{,}{/ˈevribɒdi/}{pron.}{每个人} \RV{}{even}{} people who are \RV{}{well}{} \RVB{}{trained}{}{train} \RV{}{in}{} \AG{}{maths}{,} \PH{gets\ it\ \NW{}{wrong}{}{/rɒŋ/}{adj.}{错的}}{get sth wrong}{弄错} at first, says \GL{}{psychologist}{}{/saɪˈkɒlədʒɪst/}{n.[C]}{心理学家} \PN{}{Walter Herbranson}{.}{/ˌwɔːltə ˈhɜːbrənsn/}{n.}{沃尔特·赫布兰森} To him, that only \RVB{}{shows}{}{show} their \RVB{}{brains}{}{brain} work \NS{}{normally}{.}{新义}{}{/ˈnɔːməli/}{adv.}{正常地}

\PARA{6}Look at \RV{}{each}{} of the three \NS{}{cases}{.}{新义}{case}{/keɪs/}{n.[C]}{情形} \NS{}{Say}{}{新义}{}{/seɪ/}{vt.}{假如；比如说} you \AG{}{pick}{} \AG{}{door}{} 1. If the \AG{}{car}{} is behind \AG{}{door}{} 2 or 3, the \AG{}{host}{} \RV{}{must}{} open the other goat \AG{}{door}{,} and switching \AGB{}{wins}{.}{win} Switching \RVB{}{loses}{}{lose} \PH{only\ if}{only if}{只有……才} the \AG{}{car}{} is behind \AG{}{door}{} 1.

\PARA{7}\RV{}{Yet}{} people \AG{}{keep}{} their first \AG{}{door}{} 85 to 90 \NW{}{per cent}{}{/pə ˈsent/}{n.[C]}{百分之……} of the time. \SMK{g}Psychologist \PN{}{John Petrocelli}{}{/ˌdʒɒn ˌpetrəˈtʃeli/}{n.}{约翰·彼得罗切利} \PH{had\ people\ \NW{}{repeat}{}{/rɪˈpiːt/}{vt.}{重复}}{have sb do}{让某人做} the \NS{}{game}{;}{新义}{}{/ɡeɪm/}{n.[C]}{游戏} it took many \NWB{}{rounds}{}{round}{/raʊnd/}{n.[C]}{一轮；一局} before they \AGB{}{learned}{}{learn} to switch.\EMK{g}{1} 
\GRAM{1}{\GITEM{it 表示时间}{it 可以说时间、天气、距离，本身没有具体意思。}{it took many rounds 的 it 不指东西，只说花了很多轮。}}{}
He \RVB{}{found}{}{find} out that \RV{}{part}{} of the \RV{}{reason}{} is \PH{\RV{“}{kick}{-}in-the-\NW{}{pants}{”}{/pænts/}{n.[pl.]}{裤子}}{kick-in-the-pants}{懊悔自责的} \NW{}{thinking}{:}{/ˈθɪŋkɪŋ/}{n.[U]}{思维；想法} after a \NW{}{choice}{,}{/tʃɔɪs/}{n.[C]}{选择} we think about the \AG{}{choice}{} we didn’t make. \SMK{y}\AGB{}{Staying}{}{stay} and \AGB{}{losing}{}{lose} doesn’t \NS{}{hurt}{}{新义}{}{/hɜːt/}{vi.}{让人难受} much, but if you switch and \AG{}{lose}{,} you \RV{}{feel}{} much \NWB{}{sorrier}{:}{sorry}{/ˈsɒri/}{adj.}{更懊悔的} a \RV{}{win}{} that was \NS{}{once}{}{新义}{}{/wʌns/}{adv.}{曾经} \RV{}{yours}{} is now \RVB{}{lost}{,}{lose} and that \RVB{}{feels}{}{feel} worse \RV{}{than}{} \RV{}{never}{} having it.\EMK{y}{2} 
\YI{2}{不换门而输了，不怎么难受；可要是换了门又输了，你会懊悔得多：本来属于你的胜利现在丢了，这比从来没得到过更难受。}
These \NWB{}{alternatives}{}{alternative}{/ɔːlˈtɜːnətɪv/}{n.[C]}{另一种可能} to the \RV{}{real}{} \NS{}{events}{}{新义}{event}{/ɪˈvent/}{n.[C]}{事件} can \NS{}{cloud}{}{新词性}{}{/klaʊd/}{vt.}{使模糊} your \NS{}{memory}{.}{新义}{}{/ˈmeməri/}{n.[C]}{记忆}

\PARA{8}\AG{}{Even}{} the \AG{}{famous}{} \GL{}{mathematician}{}{/ˌmæθəməˈtɪʃn/}{n.[C]}{数学家} \PN{}{Paul Erdős}{}{/ˌpɔːl ˈeədɜːʃ/}{n.}{保罗·埃尔德什} was \NWB{}{fooled}{.}{fool}{/fuːl/}{vt.}{骗；难住} When his old friend \PN{}{Andrew Vázsonyi}{}{/ˌændruː ˈvɑːʒəni/}{n.}{安德鲁·瓦若尼} \AGB{}{showed}{}{show} him the puzzle, Erdős asked for a \AG{}{reason}{} to switch. When Vázsonyi didn’t \NW{}{provide}{}{/prəˈvaɪd/}{vt.}{提供} one, Erdős asked, \PH{“What’s\ the\ \NS{}{matter}{}{新义}{}{/ˈmætə(r)/}{n.[sing.]}{问题；毛病}\ with\ you?”}{What's the matter (with sb)?}{……怎么了？} Vázsonyi \RV{}{later}{} wrote in an \NW{}{article}{,}{/ˈɑːtɪkl/}{n.[C]}{文章} “I said that I was \NS{}{sorry}{,}{新义}{}{/ˈsɒri/}{adj.}{抱歉的} but I didn’t have a \PH{\NW{}{common}{-}{/ˈkɒmən/}{adj.}{普通的}\NW{}{sense}{}{/sens/}{n.[U]}{判断力}}{common-sense}{常识性的} \NW{}{explanation}{.”}{/ˌekspləˈneɪʃn/}{n.[C]}{解释}

\PARA{9}\RV{}{Imagine}{} 100 \AG{}{doors}{,} \RV{}{with}{} a \AG{}{car}{} behind one and goats behind the other 99. You \AG{}{pick}{} \NS{}{number}{}{新义}{}{/ˈnʌmbə(r)/}{n.[C]}{号码；……号} 1. Monty opens 98 goat \AG{}{doors}{} and \NWB{}{leaves}{}{leave}{/liːv/}{vt.}{留下} only \AG{}{number}{} 72 \AG{}{closed}{.} Do you want to \AG{}{keep}{} \AG{}{number}{} 1? What about \AG{}{number}{} 72? Here you \AG{}{should}{} switch.

\PARA{10}I \RVB{}{agreed}{}{agree} that switching was \RV{}{right}{,} but I \RV{}{still}{} \NS{}{wondered}{:}{新词性}{wonder}{/ˈwʌndə(r)/}{vt.}{想知道} why \AG{}{should}{} my first \AG{}{choice}{} \NW{}{affect}{}{/əˈfekt/}{vt.}{影响} my \RV{}{second}{} \AG{}{choice}{?} “That \NS{}{idea}{}{新义}{}{/aɪˈdɪə/}{n.[C]}{想法；看法} is so \NW{}{natural}{,}{/ˈnætʃrəl/}{adj.}{自然的} but it’s \NW{}{actually}{}{/ˈæktʃuəli/}{adv.}{实际上} \AG{}{wrong}{,”} says mathematician \PN{}{Jason Rosenhouse}{.}{/ˌdʒeɪsn ˈrəʊznhaʊs/}{n.}{杰森·罗森豪斯} Monty knows where the \AG{}{car}{} is and only \RV{}{ever}{} opens a goat \AG{}{door}{.} To \RV{}{understand}{} this, you \RV{}{need}{} to \PH{put\ \RV{}{yourself}{}\ in\ Monty’s\ \NWB{}{shoes}{.}{shoe}{/ʃuː/}{n.[C]}{鞋}}{put oneself in sb's shoes}{设身处地}

\PARA{11}Last year, a friend, her father and I \PH{got\ \RV{}{together}{}}{get together}{聚在一起} to play the \AG{}{game}{.} I \NWB{}{hid}{}{hide}{/haɪd/}{vt.}{藏} a \NW{}{sweet}{,}{/swiːt/}{n.[C]}{糖果} a \NW{}{chocolate}{}{/ˈtʃɒklət/}{n.[C]}{巧克力糖} with a \NW{}{nut}{}{/nʌt/}{n.[C]}{坚果} \RV{}{inside}{,} \NS{}{under}{}{新义}{}{/ˈʌndə(r)/}{prep.}{在……下面} one of three \NS{}{cups}{}{新义}{cup}{/kʌp/}{n.[C]}{杯子} and played Monty. \PH{In\ the\ \RV{}{end}{,}}{in the end}{最后；终于} after 10 years of \PH{a\ love-\NW{}{hate}{}{/heɪt/}{n.[U]}{恨}\ \RV{}{relationship}{}}{a love-hate relationship}{又爱又恨的关系} with the puzzle, \RV{}{everything}{} \PH{made\ \NS{}{sense}{.}{新义}{}{/sens/}{n.[U]}{意义；道理}}{make sense}{讲得通}

\PARA{12}Playing the \AG{}{host}{} helps many people \AG{}{understand}{} the puzzle. Petrocelli’s \NS{}{team}{}{新义}{}{/tiːm/}{n.[C]}{团队；小组} \NS{}{found}{}{新义}{find}{/faɪnd/}{vt.}{发现} this when they \NWB{}{programmed}{}{program}{/ˈprəʊɡræm/}{vt.}{编程序} the \AG{}{game}{} for their \NS{}{study}{.}{新词性}{}{/ˈstʌdi/}{n.[C]}{研究} \RV{}{As}{} the \AG{}{host}{,} you see that you have no \NS{}{magic}{}{新词性}{}{/ˈmædʒɪk/}{adj.}{有魔力的} \RV{}{power}{} and \RV{}{almost}{} no \AGB{}{choices}{.}{choice} \PH{On\ \NS{}{average}{,}{新词性}{}{/ˈævərɪdʒ/}{n.[U]}{平均水平}}{on average}{平均来说} one \AG{}{round}{} \AG{}{in}{} three, the \NS{}{player}{}{新义}{}{/ˈpleɪə(r)/}{n.[C]}{参赛者；玩家} \AGB{}{picks}{}{pick} the \AG{}{car}{} first, and you can open \NS{}{either}{}{新词性}{}{/ˈaɪðə(r)/}{pron.}{（两者中）任一} of the other \AG{}{doors}{.} In the other two \AG{}{rounds}{,} you can open only one \AG{}{door}{.} Being the \AG{}{host}{} is \AG{}{almost}{} \NW{}{boring}{.}{/ˈbɔːrɪŋ/}{adj.}{无聊的} You only tell the \AG{}{player}{} what \NS{}{they}{}{新义}{}{/ðeɪ/}{pron.}{他或她（单数）} \RV{}{already}{} know: \NW{“}{Hey}{,}{/heɪ/}{int.}{嘿} one of these two \AG{}{doors}{} doesn’t have the \AG{}{car}{.”} You are giving the \AG{}{player}{} a \NS{}{chance}{}{新义}{}{/tʃɑːns/}{n.[C]}{机会} to get \PH{two\ \AG{}{doors}{}\ for\ the\ \RV{}{price}{}\ of\ one.}{two for the price of one}{一份价钱得两份}

\PARA{13}Now be the \AG{}{player}{} again and \AG{}{pick}{} \AG{}{door}{} 3. The \AG{}{host}{} asks if you \RV{}{would}{} like to \NS{}{change}{}{新义}{}{/tʃeɪndʒ/}{vt.}{换；调换} it for \RV{}{both}{} of the other \AG{}{doors}{.} \NW{}{Yeah}{,}{/jeə/}{int.}{是啊} \PH{of\ \NW{}{course}{}{/kɔːs/}{n.[C]}{（of course 当然）}}{of course}{当然} you \AG{}{would}{!} Then he opens \AG{}{door}{} 1: a goat. You \AG{}{already}{} knew one of your two \AG{}{doors}{} \AG{}{hid}{} a goat. “The \RV{}{situation}{} is the \RV{}{same}{} \RV{}{as}{} in the \NW{}{beginning}{,”}{/bɪˈɡɪnɪŋ/}{n.[sing.]}{开始} says psychologist \PN{}{Elisabet Tubau}{.}{/ˌelɪzəbet tuːˈbaʊ/}{n.}{伊丽莎白·图鲍}

\PARA{14}“It’s \AG{}{almost}{} like \RV{}{magic}{,”} says \GL{}{philosopher}{}{/fəˈlɒsəfə(r)/}{n.[C]}{哲学家} \PN{}{Christopher Pynes}{.}{/ˌkrɪstəfə ˈpaɪnz/}{n.}{克里斯托弗·派恩斯} \SMK{y}You are \RV{}{really}{} \NWB{}{guessing}{}{guess}{/ɡes/}{vt.}{猜} \NW{}{whether}{}{/ˈweðə(r)/}{conj.}{是否} your first \NS{}{pick}{}{新词性}{}{/pɪk/}{n.[C]}{选择} was \AG{}{right}{,} and thanks to the \AG{}{door}{} the \AG{}{host}{} opens, the \AGB{}{chances}{}{chance} of two \AG{}{doors}{} now \NS{}{rest}{}{新义}{}{/rest/}{vi.}{落在；集中于} on a \RV{}{single}{} \AG{}{door}{.}\EMK{y}{3}
\YI{3}{你其实是在猜第一次选得对不对；而多亏主持人打开的那扇门，两扇门的可能性现在都集中到了一扇门上。}

\PARA{15}In a \RV{}{talk}{} with Pynes, I \NW{}{finally}{}{/ˈfaɪnəli/}{adv.}{终于} saw my \NW{}{mistake}{.}{/mɪˈsteɪk/}{n.[C]}{错误} If you \RV{}{turn}{} on the \NS{}{TV}{}{新义}{}{/ˌtiː ˈviː/}{n.[C]}{电视机} \PH{in\ the\ \NW{}{middle}{}{/ˈmɪdl/}{n.[sing.]}{中间}\ of}{in the middle of}{在……中间} the \AG{}{show}{} and see just two \AG{}{doors}{,} \AGB{}{picking}{}{pick} one is \RV{}{truly}{} 50-50. But the \AG{}{player}{} \AG{}{still}{} \AGB{}{wins}{}{win} two times \AG{}{in}{} three by switching. \SMK{y}\AGB{}{Chances}{}{chance} \RV{}{belong}{} to the \RV{}{person}{} \RVB{}{watching}{}{watch} and to what that \AG{}{person}{} knows, not to the \AG{}{doors}{.}\EMK{y}{4}
\YI{4}{可能性取决于观看的人和他所知道的情况，而不取决于门本身。}

\PARA{16}\SMK{g}\PH{\NW{}{Above}{}{/əˈbʌv/}{prep.}{在……之上}\ all,}{above all}{最重要的是} says vos Savant \NW{}{herself}{,}{/hɜːˈself/}{pron.}{她本人} the puzzle \AGB{}{shows}{}{show} how \RV{}{hard}{} it is for us to \RV{}{believe}{} that we \RV{}{could}{} be \AG{}{wrong}{,} \RV{}{especially}{} when we \AG{}{feel}{} \RV{}{sure}{.}\EMK{g}{2}
\GRAM{2}{\GITEM{for + 人 + to do}{for + 人 说明这件事对谁来说怎样、由谁来做。}{for us 说明是谁难以相信：是我们自己。}}{}
\end{lessonbody}
\sourceline{Adapted from: “This Infamous Probability Puzzle Fools Nearly Everyone. Here’s How to Beat It”, Scientific American, 2024-07-16.}
\SelfTest{111}{\TW{classroom}
\TW{subject}
\TW{teacher}
\TW{present}
\TW{door}
\TW{closed}
\TW{should}
\TW{lesson}
\TW{explain}
\TW{magazine}
\TW{disagree}
\TW{angry}
\TW{letter}
\TW{degree}
\TW{everybody}
\TW{wrong}
\TW{per cent}
\TW{repeat}
\TW{round}
\TW{pants}
\TW{thinking}
\TW{choice}
\TW{sorry}
\TW{alternative}
\TW{fool}
\TW{provide}
\TW{article}
\TW{common}
\TW{sense}
\TW{explanation}
\TW{leave}
\TW{affect}
\TW{natural}
\TW{actually}
\TW{shoe}
\TW{hide}
\TW{sweet}
\TW{chocolate}
\TW{nut}
\TW{hate}
\TW{program}
\TW{boring}
\TW{hey}
\TW{yeah}
\TW{course}
\TW{beginning}
\TW{guess}
\TW{whether}
\TW{finally}
\TW{mistake}
\TW{middle}
\TW{above}
\TW{herself}
\TW{meet\,\nstag{新义}}
\TW{problem\,\nstag{新义}}
\TW{host\,\nstag{新义}}
\TW{show\,\nstag{新词性}}
\TW{pick\,\nstag{新义}}
\TW{double\,\nstag{新词性}}
\TW{chance\,\nstag{新词性}}
\TW{class\,\nstag{新义}}
\TW{doctor\,\nstag{新义}}
\TW{normally\,\nstag{新义}}
\TW{case\,\nstag{新义}}
\TW{say\,\nstag{新义}}
\TW{game\,\nstag{新义}}
\TW{hurt\,\nstag{新义}}
\TW{once\,\nstag{新义}}
\TW{event\,\nstag{新义}}
\TW{cloud\,\nstag{新词性}}
\TW{memory\,\nstag{新义}}
\TW{matter\,\nstag{新义}}
\TW{sorry\,\nstag{新义}}
\TW{number\,\nstag{新义}}
\TW{wonder\,\nstag{新词性}}
\TW{idea\,\nstag{新义}}
\TW{under\,\nstag{新义}}
\TW{cup\,\nstag{新义}}
\TW{sense\,\nstag{新义}}
\TW{team\,\nstag{新义}}
\TW{find\,\nstag{新义}}
\TW{study\,\nstag{新词性}}
\TW{magic\,\nstag{新词性}}
\TW{average\,\nstag{新词性}}
\TW{player\,\nstag{新义}}
\TW{either\,\nstag{新词性}}
\TW{they\,\nstag{新义}}
\TW{change\,\nstag{新义}}
\TW{rest\,\nstag{新义}}
\TW{TV\,\nstag{新义}}
\TW{stay with sb}
\TW{50-50}
\TW{… out of …}
\TW{in class}
\TW{pretty much}
\TW{get sth wrong}
\TW{only if}
\TW{have sb do}
\TW{kick-in-the-pants}
\TW{What's the matter (with sb)?}
\TW{common-sense}
\TW{put oneself in sb's shoes}
\TW{get together}
\TW{in the end}
\TW{a love-hate relationship}
\TW{make sense}
\TW{on average}
\TW{two for the price of one}
\TW{of course}
\TW{in the middle of}
\TW{above all}}
\par\write-1{LESSONEND 9 \thepage}
